\DocumentMetadata{
  pdfversion=2.0,
  pdfstandard=ua-2,
  lang=en-US,
  tagging=on,
  tagging-setup={math/setup=mathml-SE,tabsorder=structure}
}
\documentclass[11pt,letterpaper]{article}
\usepackage[margin=0.68in,includefoot]{geometry}
\usepackage{newcomputermodern}
\usepackage{amsmath,amscd,mathtools}
\usepackage{tikz,adjustbox}
\providecommand{\square}{\mdlgwhtsquare}
\usepackage[hidelinks]{hyperref}
\hypersetup{
  pdftitle={Algebra PhD Exam, May 2025},
  pdfauthor={University of Florida Department of Mathematics},
  pdfsubject={Graduate mathematics examination},
  pdfdisplaydoctitle=true,
  bookmarksopen=true
}
\setcounter{secnumdepth}{0}
\setlength{\parindent}{0pt}
\setlength{\parskip}{0.28em}
\setlength{\emergencystretch}{3em}
\setlength{\leftmargini}{2.15em}
\setlength{\leftmarginii}{2.65em}
\setlength{\leftmarginiii}{3em}
\pagestyle{empty}
\raggedbottom
\allowdisplaybreaks
\NewTaggingSocketPlug{block/list/label}{examproblem}{%
  \tagstructbegin{tag=Lbl}\tagstructend%
  \tagstructbegin{tag=\LBody}%
  \tagstructbegin{tag=\ProblemHeadingTag,title-o={\ProblemHeadingText}}%
  \tagmcbegin{tag=\ProblemHeadingTag,actualtext={\ProblemHeadingText}}%
  #2%
  \tagmcend%
  \tagstructend%
}
\newenvironment{examproblems}[2]{%
  \def\ProblemHeadingTag{#2}%
  \AssignTaggingSocketPlug{block/list/label}{examproblem}%
  \begin{enumerate}%
  \setcounter{enumi}{#1}%
  \renewcommand{\labelenumi}{\textbf{\theenumi.}}%
  \setlength{\topsep}{0.35em}%
  \setlength{\partopsep}{0pt}%
  \setlength{\itemsep}{0.48em}%
  \setlength{\parsep}{0.18em}%
}{\end{enumerate}}
\newenvironment{examsubparts}{%
  \AssignTaggingSocketPlug{block/list/label}{default}%
  \begin{enumerate}%
  \setlength{\topsep}{0.18em}%
  \setlength{\partopsep}{0pt}%
  \setlength{\itemsep}{0.16em}%
  \setlength{\parsep}{0.1em}%
}{\end{enumerate}}
\newenvironment{examinstructions}{%
  \par\begingroup\itshape
}{%
  \par\endgroup\vspace{0.18em}
}
\makeatletter
\renewcommand\subsection{\@startsection{subsection}{2}{0pt}%
  {-1.25ex plus -.25ex minus -.1ex}%
  {0.55ex plus .1ex}%
  {\normalfont\large\bfseries}}
\renewcommand{\@maketitle}{%
  \newpage\null\vskip -1em%
  {\footnotesize\bfseries
    \href{https://www.ufl.edu/}{University of Florida}, \href{https://math.ufl.edu/}{Department of Mathematics}\par}%
  \vskip -0.5em%
  \tagstructbegin{tag=H1,title={Algebra PhD Exam, May 2025}}%
  \tagpdfparaOff%
  \tagmcbegin{tag=H1}%
    {\LARGE\bfseries\href{https://gma.math.ufl.edu/exams/algebra-phd/}{Algebra PhD Exam}, May 2025\par}%
  \tagmcend%
  \tagpdfparaOn%
  \tagstructend%
  \vskip 0.25em%
  \tagmcbegin{artifact}%
    \hrule height 1pt%
  \tagmcend%
  \vskip 1em%
}
\makeatother
\title{Algebra PhD Exam, May 2025}
\author{}
\date{}
\begin{document}
\maketitle
\thispagestyle{empty}
\begin{examinstructions}
Answer seven problems. (If you turn in more, the first seven will be graded.) Within reason, you may quote theorems as long as you state them clearly.
\end{examinstructions}
\begin{examinstructions}
Note. Below ring means associative ring with identity, and module means unital module unless otherwise specified.
\end{examinstructions}
\begin{examproblems}{0}{H2}
\def\ProblemHeadingText{Problem 1.}
\item
Let \(q>1\) be a power of a prime~\(p\). Prove that there exists some field~\(F\) with \(|F|=q\).
\def\ProblemHeadingText{Problem 2.}
\item
Determine the Galois group of \(x^4-5\) over~\(\mathbb{Q}\), over \(\mathbb{Q}(\sqrt{5})\), and over \(\mathbb{Q}(\sqrt{5}\,i)\). Justify your answers.
\def\ProblemHeadingText{Problem 3.}
\item
Let~\(R\) be an integral domain, let~\(D\) be a divisible~\(R\)-module, and let~\(T\) be a torsion~\(R\)-module. Prove that \(D\otimes_R T=\{0\}\).
\def\ProblemHeadingText{Problem 4.}
\item
Prove that the group~\(G\) which is described below using generators and relations is not solvable.
\[
G=\langle x,y\mid x^6=y^2=1\rangle
\]
\def\ProblemHeadingText{Problem 5.}
\item
Let~\(R\) be a ring with identity.
\begin{examsubparts}
\item[\textnormal{(a)}]
Define what it means for a left~\(R\)-module to be projective.
\item[\textnormal{(b)}]
Show that a free~\(R\)-module is projective.
\item[\textnormal{(c)}]
Give an example of a projective~\(R\)-module that is not free.
\end{examsubparts}
\def\ProblemHeadingText{Problem 6.}
\item
Let~\(K\) be a field. Show that any vector space over~\(K\) has a basis (do not assume the vector space has finite dimension).
\def\ProblemHeadingText{Problem 7.}
\item
Prove: If a Dedekind domain has only a finite number of nonzero prime ideals then it is a principal ideal domain. (Hint: Prove first that each prime ideal is principal, use the Chinese Remainder Theorem).
\def\ProblemHeadingText{Problem 8.}
\item
Let~\(R\) be a integral domain.
\begin{examsubparts}
\item[\textnormal{(a)}]
Define what it means for a fractional ideal of~\(R\) to be invertible.
\item[\textnormal{(b)}]
Show that an invertible ideal is a projective~\(R\)-module.
\end{examsubparts}
\def\ProblemHeadingText{Problem 9.}
\item
Suppose \(R\subseteq S\) are integral domains and let \(s\in S\). Prove that the following are equivalent:
\begin{examsubparts}
\item[\textnormal{1.}]
\(s\) is integral over~\(R\);
\item[\textnormal{2.}]
\(R[s]\) is finitely generated as an~\(R\)-module;
\item[\textnormal{3.}]
There exists some ring~\(T\) such that \(R\subseteq T\subseteq S\), and \(s\in T\), and~\(T\) is finitely generated as an~\(R\)-module.
\end{examsubparts}
\def\ProblemHeadingText{Problem 10.}
\item
Suppose~\(R\) is a ring with identity (not necessarily commutative), \(M\) is a right~\(R\)-module and
\[
N'\longrightarrow N\longrightarrow N''\longrightarrow 0
\]
 is an exact sequence of left~\(R\)-modules. Show that
\[
M\otimes_R N'\longrightarrow M\otimes_R N\longrightarrow M\otimes_R N''\longrightarrow 0
\]
 is an exact sequence of abelian groups.
\def\ProblemHeadingText{Problem 11.}
\item
This problem has two parts.
\begin{examsubparts}
\item[\textnormal{(a)}]
State the Jacobson density theorem.
\item[\textnormal{(b)}]
Use it to show that if~\(K\) is a field and~\(A\) is a simple~\(K\)-algebra of finite dimension then there is a division~\(K\)-algebra~\(D\) and a finite dimensional vector space~\(V\) over~\(D\) such that \(A\simeq\operatorname{End}_D(V)\).
\end{examsubparts}
\end{examproblems}
\end{document}
